\clearpage
\section{Local geometry of coarse observations}\label{sec:geometry183}
The preceding critical inverse-moment theorem constructs a global spectral
observation. It does not construct a spatial partition. We now consolidate
four complementary statements: an abstract limitation of spectral data,
a local coexact point-support obstruction, a positive rectangular construction
for a restricted coefficient class, and a limitation of minimum-coefficient
certificates even with sharp unweighted constants. The original unrestricted
critical rectangular question remains open.

Throughout the actual lattice statements, $D\geq3$, $p=D/2$,
$X=(\mathbb Z/n\mathbb Z)^D$, $N=n^D$ and $h=D(N_c^2-1)$.
The energy is the componentwise start-weighted form
\[
 \E_a(U)=\sum_{x,i,j}a(x)|D_iU_j(x)|^2,\qquad
 D_iU_j(x)=U_j(x+e_i)-U_j(x).
\]
We do not silently substitute a variable-coefficient weighted curl form.
For a rectangular partition $\mathcal P$, the actual normalized mean map is
\[
 C_{\mathcal P}(U)_{R,j}=\sqrt{|R|}\operatorname{mean}_R U_j,
 \qquad \dim F=h\#\mathcal P.
\]
A global inequality with penalty $c\|C_{\mathcal P}(U)\|^2$ gives the
original differentiable near-zero data with $K=c$. A kernel obstruction
below applies to every finite locally upper-quadratic penalty, rather
than only to this particular coefficient.

\subsection{An isospectral obstruction to forgetting coordinate geometry}
This first construction is deliberately a nonlocal scalar graph, not an
admissible nearest-neighbour Hodge profile. Let $n$ be even and set
$e=N^{-1/2}\mathbf1$, $v=N^{-1/2}(-1)^{x_1+\cdots+x_D}$. For $0<\alpha<c<\beta$ put
\[
 L=\beta(I-ee^*)-(\beta-\alpha)vv^*.
\]
Its eigenvalues are $0$ on $e$, $\alpha$ on $v$, and $\beta$ on their
orthogonal complement. On $e^\perp$ there is exactly one mode below $c$,
and the actual scalar observation $\langle v,f\rangle$ with penalty
coefficient $c-\alpha$ repairs the target. For $x\ne y$ the edge weight is
\[
 w_{xy}=\frac{\beta+(\beta-\alpha)\sigma(x)\sigma(y)}N>0,
 \qquad \sigma(x)=(-1)^{x_1+\cdots+x_D}.
\]
The row sums vanish, so this is a connected positive graph Laplacian.
For $\alpha=c/2$, $\beta=2c$, the weight ratio is only seven.

Let $\Pi$ project onto the span of normalized cell indicators and let
$\tau=\|\Pi v\|^2$. Constants belong to this span. If $\tau<1$, the actual
unseen mean-zero vector $u=(I-\Pi)v$ has
\[
 \|u\|^2=1-\tau,\qquad
 \frac{\langle Lu,u\rangle}{\|u\|^2}
       =\alpha+(\beta-\alpha)\tau.
\]
For a rectangle, the alternating sum in every cyclic coordinate interval
has magnitude at most one. Consequently, for $q$ cells,
\[
 \tau=\frac1N\sum_R\frac{(\sum_R\sigma)^2}{|R|}\leq\frac qN.
\]
The exact finite quadratic-penalty criterion is stricter than a nonnegative
kernel margin. Put $d=\beta-c$ and $r=\beta-\alpha$. On $e^\perp$,
$L-cI+K\Pi=dI+K\Pi-rvv^*$. Since $dI+K\Pi>0$, Cauchy--Schwarz after
its square root shows that positivity is equivalent to
\[
 r\left[\frac{1-\tau}{d}+\frac\tau{d+K}\right]\leq1.
\]
Solving gives
\[
 \text{some finite }K\text{ exists}\quad\Longleftrightarrow\quad
 r\tau>c-\alpha,
 \qquad
 K_{\min}=\frac{(\beta-c)(c-\alpha)}{(\beta-\alpha)\tau-(c-\alpha)}.
\]
At equality the second summand stays positive for every finite $K$.
A locally upper-quadratic penalty cannot outperform this quadratic
condition. Thus the displayed positive-weight example requires $q>N/3$
for every successful rectangular mean partition. Its restricted mean-zero
observation rank is $q-1$.

Replace the checkerboard signs by signs constant on two equal half-slabs.
A vertex permutation conjugates the two graphs: their complete spectra
and unlabelled weighted graphs coincide. The two half-slab means observe
the sole low direction and admit a finite penalty. Therefore even the
unlabelled graph, not just the eigenvalue count, omits the coordinate
geometry needed for rectangular observations. This is an analytical
counterexample to a locality-free conversion; it is not a counterexample
to the original local coefficient question or a new Schur principle.

\subsection{Growing point support on actual coexact fields in dimension three}
The earlier line-constant point test can be strengthened using a local
coexact correction, with an essential change of coefficient scale.
For even $n\geq4$, define the rational triangular profile and its periodic
primitive by
\[
 f(t)=1-4\operatorname{dist}_{\mathbb Z/n\mathbb Z}(t,0)/n,
 \qquad F(0)=0,\quad F(t)=-\sum_{r=1}^t f(r).
\]
Then $\sum f=0$, $|D f|=4/n$, $D^*F=f$ and
$\sum f^2=n/3+8/(3n)$. The positive-pair curl and backward adjoint satisfy
\[
 (d_1U)_{ij}=D_iU_j-D_jU_i,\qquad
 (d_1^*G)_1=-D_2^*G_{12},\quad (d_1^*G)_2=D_1^*G_{12}.
\]
Put $B^0_{12}(x)=F(x_1)e_0$ in one fixed internal unit coordinate, with
other entries zero. Then $U^0=d_1^*B^0$ has only component two, equal to
$f(x_1)e_0$. It is explicitly coexact and divergence zero, and
\[
 \|U^0\|^2\geq N/3,\qquad \E_\mu(U^0)=16\mu N/n^2.
\]

For an arbitrary sampled set $S$ of $s$ vertices, join two samples when
their periodic maximum-coordinate distance is at most four. Let $S_\alpha$
be the resulting components, of sizes $s_\alpha$, with representatives
$p_\alpha$. Write $N_r(S_\alpha)$ for the radius-$r$ neighbourhood and set
\[
 \chi_\alpha=\mathbf1_{N_1(S_\alpha)},\qquad
 (G_\alpha)_{12}(x)=\chi_\alpha(x)
       [F(x_1)-F((p_\alpha)_1)]e_0,
 \qquad U=U^0-\sum_\alpha d_1^*G_\alpha.
\]
This is an actual potential, not an assumed right inverse of divergence.
At every sample, $\chi_\alpha$ equals one at the site and its two backward
neighbours, so $d_1^*G_\alpha$ agrees with $U^0$ in component two and is
zero in component one. Other components contribute zero there. Thus
$U_i(x)=0$ for every $x\in S$ and every $i$.

Distinct $N_2(S_\alpha)$ are disjoint and their total size is at most
$5^Ds$. Outside them $U=U^0$, hence, with $C_D=5^D$,
\[
 \|U\|^2\geq N/3-C_Ds.
\]
A component graph path has at most $s_\alpha-1$ steps of cyclic length
at most four. The unit increment bound for $F$ gives
$|F(x_1)-F((p_\alpha)_1)|\leq4s_\alpha$ on $N_1(S_\alpha)$.
Each correction has at most two components of magnitude $8s_\alpha$,
so its squared norm, summed over disjoint supports, is at most
$128C_D\sum_\alpha s_\alpha^3$. Expanding finite differences gives the
unit flat Hodge identity, and $\E_\mu(A)\leq4D\mu\|A\|^2$.
The quadratic inequality for $U^0$ minus the correction therefore gives
\[
 \E_\mu(U)\leq32\mu N/n^2+A_D\mu\sum_\alpha s_\alpha^3,
 \qquad A_D=1024D5^D.                                      \tag{G1}
\]
All support and energy constants have been derived, without a separation
or correction-operator hypothesis on the samples.

\begin{proposition}\label{prop:geo-point}
Fix $c>0$ and take $a=\mu_n=cn^2/384$ for sufficiently large even $n$.
Any successful point-supported original observation on the actual coexact
domain must have sampled support
\[
 s\geq\min\left\{\frac{n^D}{12C_D},
       (64/A_D)^{1/3}n^{(D-2)/3}\right\}.
\]
The same conclusion holds when only the derivative at zero factors through
the point values. In particular support diverges in $D=3$.
\end{proposition}
\begin{proof}
The profile satisfies $1\leq a\leq2\mu_n$, its mean is $\mu_n$, and its
normalized critical inverse moment is exactly one. If $s\leq N/(12C_D)$,
then $\|U\|^2\geq N/4>0$. Point-supported maps with $C(0)=0$ annihilate
$tU$; derivative-supported maps satisfy $C(tU)=o(t)$. Divide the local
coercivity inequality by $t^2$ and use the quadratic upper bound on its
penalty. Since $32\mu_nN/n^2=cN/12$, (G1) implies
$\sum s_\alpha^3\geq(64/A_D)n^{D-2}$. Finally $\sum s_\alpha^3\leq s^3$.
If the initial restriction on $s$ fails, the first alternative already holds.
\end{proof}
The family scale is $\mu_n=cn^2/384$, not the older $\mu=n^2,c\leq16$
coexact regime, where a uniform Fourier gap can already suffice without
observations. No fine-space repair is covered here. On this same new
family, observing all real Fourier modes with centred $|k|^2<24$ has
dimension at most $h9^D$ and gives target $c$ with penalty coefficient $c$:
the omitted modes obey $\mu_n\lambda_n(k)\geq c|k|^2/24$.
Thus this separates bounded point support from bounded global spectral
rank, not from rectangular averages or point samples plus complementary
means. The arbitrary-cluster exponent is a safe bound, not claimed sharp.

\subsection{A fixed-gap rectangular construction with arbitrary weak geometry}
\begin{lemma}\label{lem:geo-isolate}
Suppose $a\geq A>0$ outside any set $S$ of $k$ vertices, with $a>0$
everywhere. Put $L=\min\{n,\max\{1,\lfloor\sqrt{A/c}\rfloor\}\}$.
There is an actual rectangular partition satisfying
\[
 c\|U\|^2\leq\E_a(U)+c\|C_{\mathcal P}(U)\|^2
 \quad\hbox{for every cochain},\qquad
 \#\mathcal P\leq\lceil n/L\rceil^D+2Dk.                  \tag{G2}
\]
No shape, connectedness or separation of $S$ is assumed.
\end{lemma}
\begin{proof}
Start with coordinate intervals of length at most $L$. To isolate a sample
in its current rectangle, cut just before and after its first coordinate,
retain the side slabs, and continue only in the central width-one slice.
Repeat in each direction. At most $2D$ actual binary cuts isolate that
vertex; each increases the leaf count by one. Previously isolated vertices
stay singleton leaves. Induction gives disjoint coverage, the asserted
count, and every nonsingleton leaf avoids $S$. This is not a product of
all coordinate cut sets or a full equal-child cube tree.

On a length-$l$ path, the counting variance identity, path Cauchy--Schwarz
and the pair count $r(l-r)\leq l^2/4$ give
\[
 \sum|u-\bar u|^2=\frac1l\sum_{r<s}|u_r-u_s|^2
 \leq\frac{l(l-1)}4\sum|Du|^2\leq\frac{l^2}4\sum|Du|^2.
\]
Commuting coordinate-averaging projections $P_i$ satisfy
$I-\prod_iP_i\leq\sum_i(I-P_i)$ as quadratic forms; their simultaneous
zero-one eigenvalues prove it. Tensorization bounds each leaf's variance
by $L^2/4$ times its unweighted internal-edge energy. Every retained
positive edge starts outside $S$, hence has coefficient at least $A$.
If $L\geq2$ this bounds $c$ times residual by one quarter of internal
weighted energy; if $L=1$ every residual is zero. Sum the disjoint internal
edge sets and discard nonnegative cross-leaf/wrap edges. The exact
orthogonal mean decomposition proves (G2). Arbitrary aspect ratios cause
no extra hypothesis.
\end{proof}
A safe mesh count is
\[
 \lceil n/L\rceil^D\leq4^D\max\{1,c^p\}\max\{1,N/A^p\}.
\]
For $A/c<4$, $L=1$ and $N\leq(4c)^pN/A^p$; otherwise, unless $L=n$,
$L\geq\sqrt{A/c}/2$ and $\lceil n/L\rceil\leq1+2n\sqrt{c/A}$.
These estimates also include small sizes and all-singleton cases.

\begin{theorem}\label{thm:geo-gap}
Fix $B\geq1$ and $0<\theta\leq1$ independently of $n,\mu$. In addition
to the original critical-moment profile hypotheses, assume
\[
 \text{every }x\text{ satisfies }a(x)\leq B\text{ or }a(x)\geq\theta\mu.
                                                                    \tag{G3}
\]
Then the original rectangular inequality holds globally with $K=c$ and
\[
 \dim F\leq h2^p\left[4^D\max\{1,c^p\}\theta^{-p}+2DMB^p\right]
             \max\{1,N/(1+\mu)^p\}.
\]
\end{theorem}
\begin{proof}
For $S=\{a\leq B\}$ the moment hypothesis derives
$|S|\leq MB^pN/\mu^p$. Off $S$, (G3) gives $a\geq\theta\mu$.
Insert these facts in (G2), use $\theta\leq1$ and
$\max\{1,N/\mu^p\}\leq2^p\max\{1,N/(1+\mu)^p\}$.
\end{proof}
In particular every original profile taking values $\{1,H\}$ is included:
its actual mean gives $H\geq\mu$. Arbitrary weak-set geometry replaces the
earlier single weak cube. This does not cover every two-valued family
whose lower value grows with $n$; (G3) is additional and is not derived
from the critical moment alone.

The limitation is explicit. At $c=1$, set $n=t^2$, $\mu=t^4$, $b=t^2$,
let the weak cube have side $t$ and fraction $s=t^{-D}$, and set
$a=b$ there and $H=(\mu-sb)/(1-s)$ elsewhere. The exact mean is $\mu$,
$1\leq a\leq2\mu$, and the critical moment is less than two.
Any uniform strong mesh followed by threshold-singleton isolation costs
at least $t^D$ cells: below threshold $b$ the mesh side is at most $t$;
at or above $b$ there are at least $t^D$ mandated singletons.
Nevertheless the weak cube itself and its $D$ first-coordinate-outside
complement slabs form $D+1$ rectangles. Their variance coefficients are
at most $t^2/(4b)=1/4$ or $n^2/(4H)\leq1/4$, respectively.
They therefore give the actual target with $K=1$. This defeats the named
threshold method, not unrestricted rectangular means.

\subsection{Actual detours and the failure of minimum-only certificates}
We next leave coefficient minima behind, using the classical path
comparison mechanism\cite{geo:comparison}. Let a weak set $W$ have periodic
maximum-coordinate separation at least three, let $a\geq H>0$ off $W$,
and allow arbitrary nonnegative coefficients on $W$. Denote by
$\E_{\rm str}$ and $\E_{\rm wk}$ the unweighted edge-square sums with
positive starts outside and inside $W$.

\begin{lemma}\label{lem:geo-detour}
For $D\geq2,n\geq3$ under these hypotheses,
\[
 \E_{\rm wk}\leq3\E_{\rm str},\qquad
 \E_0\leq4\E_{\rm str},\qquad \E_a\geq(H/4)\E_0.         \tag{G4}
\]
\end{lemma}
\begin{proof}
Let $\sigma(i)=i+1$ cyclically. Bypass every weak-start edge $x\to x+e_i$
by the actual three-edge path
\[
 x\longrightarrow x-e_{\sigma(i)}
 \longrightarrow x+e_i-e_{\sigma(i)}\longrightarrow x+e_i.
\]
Its positive starts are $x-e_{\sigma(i)}$ twice in different directions,
and $x+e_i-e_{\sigma(i)}$. They are strong by separation. The first edge
is traversed backwards, so its coefficient is not evaluated at $x$.

For a fixed positive edge $(y,k)$, the three possible weak sources are
\[
 y+e_k,\qquad y+e_{\sigma(k)},\qquad
 y-e_{\sigma^{-1}(k)}+e_k.
\]
For each type its direction $i$ is uniquely fixed. These candidate sources
are distinct and pairwise at periodic maximum-coordinate distance at most
two. At most one is weak. Thus every strong positive edge is used at most
once, counting uses rather than just distinct paths. Cauchy--Schwarz along
each three-step path and summation prove $\E_{\rm wk}\leq3\E_{\rm str}$.
The remaining assertions follow directly, componentwise. Periodic wraps
are already included in the path identities.
\end{proof}
The Fourier torus gap gives
\[
 \E_a(U)\geq\frac H4\frac{16}{n^2}\|U-\operatorname{mean}_XU\|^2.
                                                                    \tag{G5}
\]
If $4H/n^2\geq c$, one global normalized mean gives the original inequality
with $K=c$, on all cochains. No supplied flow, comparison or kernel margin
remains as an unproved assumption.

For a nonsingleton rectangle $R$, let $\lambda_1^0(R)$ be the sharp
unweighted internal Neumann gap after cutting each coordinate interval
into a path. A standard sufficient local test is
\[
       \min_R(a)\lambda_1^0(R)\geq c.                      \tag{G6}
\]
Singletons have zero residual. On a longest side of length $l$, the actual
mean-zero ramp $u(r)=r-(l-1)/2$ has energy $l-1$ and counting squared norm
$l(l^2-1)/12$. Tensor copies cancel, giving
\[
       \lambda_1^0(R)\leq12/[l(l+1)]<12/l^2.              \tag{G7}
\]
Thus the following obstruction applies even to the best unweighted
constant, not just a loose path or trace estimate.

\begin{proposition}\label{prop:geo-minimum}
For every fixed $D\geq3$ and $M>1$ there are original admissible profiles
at $c=1,\mu=n^2$ for which a single rectangular mean gives original data
with $K=1$, but every partition passing (G6) has an unbounded number of
cells. No fixed multiplicative loosening of (G6) rescues its uniform rank
order. This does not refute the original rectangular question.
\end{proposition}
\begin{proof}
Choose an integer $m\geq4$ with $1+m^{-D}\leq M$, and take large multiples
$t$ of $m$. Put $n=t^2$, $\mu=t^4$, $b=t^2$ and
\[
 \begin{gathered}
 W=\{0,mt,2mt,\ldots,n-mt\}^D,\quad
 k=(t/m)^D,\quad s=\frac1{m^Dt^D},\\
 a=b\text{ on }W,\qquad H=\frac{\mu-sb}{1-s}\text{ elsewhere}.
 \end{gathered}
\]
Then $H\geq\mu$, $1\leq a\leq2\mu$, and the arithmetic mean is exactly
$\mu$. The normalized critical inverse moment is
\[
 s(\mu/b)^p+(1-s)(\mu/H)^p
      =m^{-D}+(1-s)(\mu/H)^p<1+m^{-D}\leq M.
\]
Weak sites have periodic separation $mt$. By (G5), $\E_a$ is at least
four times global mean-zero squared norm, so one mean suffices. The target
rank scale is exactly one.

Every weak-containing nonsingleton leaf passing (G6) has, by (G7),
$l^2<12b=12t^2$, hence $l<\sqrt{12}t<mt$. It contains at most one weak
site; a singleton does too. Therefore every passing partition has
$\#\mathcal P\geq k=(t/m)^D\to\infty$, with arbitrary cut locations and
translated cyclic intervals allowed. For the test
$\eta\min_R(a)\lambda_1^0(R)\geq1$ with any fixed $\eta>0$, choose also
$m>\sqrt{12\eta}$. The same proof applies.
\end{proof}
If a local test retains periodic edges on a full-period side, the ramp
bound (G7) must not be used on that cycle. Its first Fourier mode instead
gives a gap upper bound $4\pi^2/l^2<40/l^2$. Choosing
$m>\sqrt{40\eta}$ gives the same conclusion uniformly. Alternatively,
for every fixed $m$, a weak-containing full-period side fails that test
for sufficiently large $t$, since its upper bound is
$4\eta\pi^2/t^2<1$. No periodic edge was discarded in the actual global
comparison (G4)--(G5).

A separate example clarifies non-necessity of the scalar moment for an
individual family. Let $n$ run through multiples of three, take
$W=\{0,3,\ldots,n-3\}^D$, $s=3^{-D}$, $a=1$ on $W$, $\mu=n^2$ and
$H=(\mu-s)/(1-s)$ elsewhere. The original mean and coefficient bounds
hold, and (G5) again gives one-cell data at $c=K=1$, while
\[
 \operatorname{mean}((\mu/a)^{D/2})\geq3^{-D}n^D\longrightarrow\infty.
\]
This does not remove the moment assumption from its unrestricted
sufficiency theorem or contradict its sharp exponent among universal
scalar inverse-moment assumptions. Here separation is a verified feature
of the explicit profiles, not a consequence of an arbitrary moment bound.

\subsection{Reproduction and remaining limitations}
The analytical sources are the pinned notes at
\texttt{aaa9e278f2c7} (nonlocal graph), \texttt{bfc4ada98dd9} (coexact
point support), \texttt{ee2b86c0e62b} (fixed-gap rectangles) and
\texttt{29fc9366bf4e} (detours and minimum certificates).
Classical comparison is attributed to Section2B and Example2.1
of\cite{geo:comparison}, which already use short bypasses and edge
congestion. No new comparison principle or independent originality is
claimed. The present finite start-weighted and coordinate-geometry
adapters are derived in the proofs above.

\begin{center}\small
\begin{tabular}{llll}
\toprule
Replay & Executed source & Cases & Rejected mutations\\
\midrule
178 & \texttt{115fedf62ad2} & 17 & 7\\
180 & \texttt{1cf122830b45} & 12 & 5 plus2 solver attacks\\
182 & \texttt{4422b0e19acf} & 11 & 7 plus2 solver attacks\\
\bottomrule
\end{tabular}
\end{center}
Each replay used four fresh pinned clones across two own CPU Colab
allocations and eight byte-identical normal/optimized payloads. Source
hashes, full stdout/stderr, actual decimal child exit codes, archive
manifests, resource records and disconnections are retained in the
reproducible package. Windows performed editing, git and hash/readback
work, not scientific execution. In180 the two host fingerprints matched;
in178 and182 they differed. Neither observation establishes independent
model review or physical-machine independence.

Replay178 checks the actual coexact potential, every pin component,
backward divergence, norm and Hodge/gradient identities. Replay180
includes ten exact $64\times64$ positivity checks and two intermediate
island cases; no large full matrix was claimed. Replay182 includes four
exact small positivity checks, the full path/preimage/edge-use dictionary,
sparse large-family coefficient checks and the path-versus-cycle ramp
exception. The nonlocal graph proof has no corresponding executed replay
in this consolidation. Finite selected cases do not prove all-size
statements, formalize these additions, establish novelty or solve the
unrestricted rectangular question.

The next substantive positive mechanism must use actual weighted local
energy or capacity rather than eigenvalue count alone, point support,
full cube trees, trace domination, a threshold-singleton strategy or
minimum coefficient alone. These scoped negative results preserve the
successful partitions on the same examples. They make no statement about
interacting gauge forms, a continuum limit or a physical mass gap.